\section{Installation}
\label{sec:install}

\paragraph{Requirements.} A C11 compiler (GCC or Clang), CMake 3.10 or
later, and OpenMP (optional: experiments in parallel). For the scripts:
Python 3 with \cmd{numpy}, \cmd{matplotlib} and \cmd{ortools} (the exact
solver only). For the paper and this manual: a LaTeX distribution (on
Debian: \cmd{texlive-latex-extra texlive-science texlive-pictures
texlive-fonts-extra}).

\paragraph{Build and test.}
\begin{shell}
cmake -S . -B build\\
cmake -{}-build build -j\\
cd build \&\& ctest -{}-output-on-failure\\
pip install -r requirements.txt \quad\# numpy, matplotlib, ortools
\end{shell}
The programs are in \cmd{build/bin}. All tests must pass
(\cref{sec:validation}); the exact-solver test runs only if OR-Tools is
installed. A build with the address and undefined-behaviour sanitizers
checks the memory accesses of the C code:
\begin{shell}
cmake -S . -B build-san -DMP\_SANITIZE=ON -DCMAKE\_BUILD\_TYPE=Debug
\end{shell}

\needspace{10\baselineskip}
\paragraph{First runs.} Each command below takes well under a second.
\begin{shell}
build/bin/mapping \quad\# a generated application, built-in zonal platform\\
build/bin/mapping examples/uc1\_central\_sizing/application.json \textbackslash\\
\quad examples/uc1\_central\_sizing/configuration.json -{}-cores\\
build/bin/mapping -{}-example-config domain > domain.json \quad\# a template\\
build/bin/mapping -{}-config domain.json -{}-chains 120 -{}-out results-domain
\end{shell}
The first run generates 300 task-chains and compares the four allocation
rules on the built-in zonal platform. The second maps the 24 functions of
use case~1 (\cref{sec:uc1}) and prints every core and every task-chain.
The last two write a template configuration file and map a generated
application on it, the results going to the files of
\cref{sec:outfiles}.
